% GENERATED FILE. Edit canonical lessons, then run npm run generate:books.
% canonical-source-hash: fnv1a64:ecaf3d78fe9bd4a2
% canonical-lessons: GE-C476-stur, GE-C476-hartnaeckig, GE-C476-launisch, GE-C476-gereizt, GE-C476-empfindlich

\chapter{Stur und launisch}
\label{ch:ge-stur-und-launisch}
\hlchaptermodality{476}

\begin{chapteropening}
\textbf{By the end of this chapter:} \emph{I can describe difficult characters.}

Complete the last lesson of chapter 476: I can describe difficult characters.
\end{chapteropening}

\section[stur]{stur — stubborn}

\label{lesson:GE-C476-stur}



\begin{quote}
Before the new one: say the German for modest, then the German for vain.
\end{quote}

\subsection*{You'll want to know: stur}
\textbf{stur} — \textquotedblleft{}stubborn\textquotedblright{}.

\begin{quote}
Mein Opa ist stur und ändert nie seine Meinung.
\end{quote}

\begin{quote}
\emph{My grandpa is stubborn and never changes his mind.}
\end{quote}

Informal: \emph{stur wie ein Esel} — stubborn as a mule.

\begin{cousinweb}
\emph{stur} is usually traced to Low German \emph{stūr} \textquotedblleft{}stiff, unyielding\textquotedblright{}, which spread into standard German; EtymDB lists the German word without a source.
\end{cousinweb}

\begin{grammarlens}[title={What you've built}]
One more word for telling, explaining and describing at B1.
\end{grammarlens}

\subsection*{Your turn}
\begin{itemize}
\raggedright
  \item \emph{Say it:} \emph{stur}
  \item \emph{Say it:} \emph{stur}, once more
  \item \emph{Say it:} a sentence of your own with \emph{stur}
  \item \emph{Recall:} say the German for modest, then the German for vain, then say \emph{stur} again
\end{itemize}

\begin{tcolorbox}[breakable,colback=teal!4,colframe=teal!35!black,title={Before you move on}]
What does stur mean? (\textquotedblleft{}Stubborn\textquotedblright{}.) Say it once more.
\end{tcolorbox}
\section[hartnäckig]{hartnäckig — persistent, dogged}

\label{lesson:GE-C476-hartnaeckig}



\begin{quote}
Before the new one: say the German for vain, then the German for stubborn.
\end{quote}

\subsection*{You'll want to know: hartnäckig}
\textbf{hartnäckig} — \textquotedblleft{}persistent, dogged\textquotedblright{}.

\begin{quote}
Der Husten ist hartnäckig und geht nicht weg.
\end{quote}

\begin{quote}
\emph{The cough is persistent and won't go away.}
\end{quote}

Of people too: \emph{Er fragte hartnäckig weiter.}

\begin{cousinweb}
\emph{hartnäckig} is built on \emph{hart} \textquotedblleft{}hard\textquotedblright{} and the word for the back of the neck: someone with a stiff neck who will not bow. \emph{Hart} is inherited from Old High German \emph{hart}, from Proto-Germanic \emph{harduz} (EtymDB).
\end{cousinweb}

\begin{grammarlens}[title={What you've built}]
One more word for telling, explaining and describing at B1.
\end{grammarlens}

\subsection*{Your turn}
\begin{itemize}
\raggedright
  \item \emph{Say it:} \emph{hartnäckig}
  \item \emph{Say it:} \emph{hartnäckig}, once more
  \item \emph{Say it:} a sentence of your own with \emph{hartnäckig}
  \item \emph{Recall:} say the German for vain, then the German for stubborn, then say \emph{hartnäckig} again
\end{itemize}

\begin{tcolorbox}[breakable,colback=teal!4,colframe=teal!35!black,title={Before you move on}]
What does hartnäckig mean? (\textquotedblleft{}Persistent, dogged\textquotedblright{}.) Say it once more.
\end{tcolorbox}
\section[launisch]{launisch — moody}

\label{lesson:GE-C476-launisch}



\begin{quote}
Before the new one: say the German for stubborn, then the German for persistent.
\end{quote}

\subsection*{You'll want to know: launisch}
\textbf{launisch} — \textquotedblleft{}moody\textquotedblright{}.

\begin{quote}
Am Morgen ist meine Schwester oft launisch.
\end{quote}

\begin{quote}
\emph{My sister is often moody in the morning.}
\end{quote}

Built on \emph{die Laune} — mood: \emph{gute Laune}, \emph{schlechte Laune}.

\begin{cousinweb}
\emph{launisch} is built on \emph{Laune} \textquotedblleft{}mood\textquotedblright{}, which goes back to Middle High German \emph{lūne} and is usually traced to Latin \emph{lūna} \textquotedblleft{}moon\textquotedblright{} (Lewis \& Short): moods were thought to change with the moon.
\end{cousinweb}

\begin{grammarlens}[title={What you've built}]
One more word for telling, explaining and describing at B1.
\end{grammarlens}

\subsection*{Your turn}
\begin{itemize}
\raggedright
  \item \emph{Say it:} \emph{launisch}
  \item \emph{Say it:} \emph{launisch}, once more
  \item \emph{Say it:} a sentence of your own with \emph{launisch}
  \item \emph{Recall:} say the German for stubborn, then the German for persistent, then say \emph{launisch} again
\end{itemize}

\begin{tcolorbox}[breakable,colback=teal!4,colframe=teal!35!black,title={Before you move on}]
What does launisch mean? (\textquotedblleft{}Moody\textquotedblright{}.) Say it once more.
\end{tcolorbox}
\section[gereizt]{gereizt — irritable, irritated}

\label{lesson:GE-C476-gereizt}



\begin{quote}
Before the new one: say the German for persistent, then the German for moody.
\end{quote}

\subsection*{You'll want to know: gereizt}
\textbf{gereizt} — \textquotedblleft{}irritable, irritated\textquotedblright{}.

\begin{quote}
Nach der langen Fahrt waren alle müde und gereizt.
\end{quote}

\begin{quote}
\emph{After the long drive everyone was tired and irritable.}
\end{quote}

The participle of \emph{reizen} — to irritate, to provoke.

\begin{cousinweb}
\emph{gereizt} is the participle of \emph{reizen} \textquotedblleft{}to provoke, to irritate\textquotedblright{}, a native verb usually traced to Old High German \emph{reizzen}; EtymDB lists it without a source.
\end{cousinweb}

\begin{grammarlens}[title={What you've built}]
One more word for telling, explaining and describing at B1.
\end{grammarlens}

\subsection*{Your turn}
\begin{itemize}
\raggedright
  \item \emph{Say it:} \emph{gereizt}
  \item \emph{Say it:} \emph{gereizt}, once more
  \item \emph{Say it:} a sentence of your own with \emph{gereizt}
  \item \emph{Recall:} say the German for persistent, then the German for moody, then say \emph{gereizt} again
\end{itemize}

\begin{tcolorbox}[breakable,colback=teal!4,colframe=teal!35!black,title={Before you move on}]
What does gereizt mean? (\textquotedblleft{}Irritable, irritated\textquotedblright{}.) Say it once more.
\end{tcolorbox}
\section[empfindlich]{empfindlich — sensitive, touchy}

\label{lesson:GE-C476-empfindlich}



\begin{quote}
Before the new one: say the German for moody, then the German for irritable.
\end{quote}

\subsection*{You'll want to know: empfindlich}
\textbf{empfindlich} — \textquotedblleft{}sensitive, touchy\textquotedblright{}.

\begin{quote}
Im Sommer ist meine Haut sehr empfindlich.
\end{quote}

\begin{quote}
\emph{In summer my skin is very sensitive.}
\end{quote}

Of people: easily hurt, touchy.

\begin{cousinweb}
\emph{empfindlich} is built on \emph{empfinden} \textquotedblleft{}to feel, to sense\textquotedblright{} with \emph{-lich}; \emph{empfinden} is \emph{finden} \textquotedblleft{}to find\textquotedblright{} with the prefix \emph{ent-} (EtymDB lists both German words).
\end{cousinweb}

\begin{grammarlens}[title={What you've built}]
One more word for telling, explaining and describing at B1.
\end{grammarlens}

\subsection*{Your turn}
\begin{itemize}
\raggedright
  \item \emph{Say it:} \emph{empfindlich}
  \item \emph{Say it:} \emph{empfindlich}, once more
  \item \emph{Say it:} a sentence of your own with \emph{empfindlich}
  \item \emph{Recall:} say the German for moody, then the German for irritable, then say \emph{empfindlich} again
\end{itemize}

\begin{tcolorbox}[breakable,colback=teal!4,colframe=teal!35!black,title={Before you move on}]
What does empfindlich mean? (\textquotedblleft{}Sensitive, touchy\textquotedblright{}.) Say it once more.
\end{tcolorbox}
